
\documentclass[12pt,a4paper]{report}

% Packages
\usepackage{amsmath}
\usepackage{graphicx}
\usepackage{braket}
\usepackage{booktabs}
\usepackage{array}

% Document information
\title{MZI (integrated optics project)}
\author{Adreet Sarkar\\
Naitik Gupta\\
Pratiksha Barman\\
Sayanti Pandit}
\date{}

\begin{document}

\maketitle

\chapter{Introduction}



\section{Material Selection}

The material parameters used in the simulation
are summarized in Table~\ref{tab:parameters}.

\begin{table}[htbp]
    \centering
    \caption{Material and optical parameters}
    \label{tab:parameters}

    \begin{tabular}{lll}
        \toprule
        \textbf{Parameter} &
        \textbf{Material} &
        \textbf{Refractive Index} \\
        \midrule

        $n_1$ (core) &
        Lithium niobate (LiNbO$_3$) &
        2.138 \\

        $n_2$ (cladding) &
        Silicon dioxide (SiO$_2$) &
        1.444 \\

        Operating wavelength &
        1550 nm &
        --- \\

        \bottomrule
    \end{tabular}
\end{table}

\section{Checking single-mode condition}

If waveguide width
\[ w=0.45\,\mu\mathrm m
\]

For an ideal symmetric slab waveguide, define
\[
V=\frac{2\pi}{\lambda}w\sqrt{n_1^2-n_2^2}.
\]

For this definition, the first higher-order mode reaches cutoff at \(V=\pi\).

Therefore, the ideal single-mode criterion is:
\[
\boxed{V<\pi}
\]

Substituting:
\[
V=\frac{2\pi}{1.55}(0.45)\sqrt{2.138^2-1.444^2}
\]

\[
\boxed{V\approx2.876<3.142}
\]

 selected values satisfy this criterion, although they are relatively close to the first higher-order mode cutoff.

\end{document}
